% Generated draft: requires mathematical and source review.
\documentclass[11pt]{article}
\usepackage[margin=25mm]{geometry}
\usepackage{fontspec}
\setmainfont{Latin Modern Roman}
\usepackage{amsmath,amssymb,mathtools,unicode-math}
\setmathfont{Latin Modern Math}
\usepackage{enumitem,needspace,xurl,hyperref}
\setlength{\parindent}{0pt}
\setlength{\parskip}{5pt}
\setlength{\emergencystretch}{3em}
\newcommand{\N}{\mathbb{N}}
\newcommand{\Z}{\mathbb{Z}}
\newcommand{\Q}{\mathbb{Q}}
\newcommand{\R}{\mathbb{R}}
\newcommand{\C}{\mathbb{C}}
\newcommand{\F}{\mathbb{F}}
\newcommand{\A}{\mathbb{A}}
\newcommand{\PP}{\mathbb{P}}
\newcommand{\E}{\mathbb{E}}
\newcommand{\eps}{\varepsilon}
\newcommand{\dd}{\,\mathrm{d}}
\begin{document}
\section*{UnsolvedMath ID: 800690}
\textbf{Source problem code:} catalogue entry
\subsection{Definitions and mathematical statement}
The supplied question asks whether a holomorphic foliation \(\mathscr{F}\), described by forms \(\alpha_j\) on \(\mathbb{D}^k\), can always be approximated by foliations on \(\C^k\), respectively \(\PP^k\). Here \(\mathbb{D}^k=\{z\in\C^k:|z_\ell|<1\text{ for }1\leq\ell\leq k\}\). The supplied excerpt does not specify the number or degrees of the forms, their rank assumptions, or the condition referenced as ``the above Frobenius condition.'' Consequently, the following is a conditional formulation under the usual interpretation by holomorphic Pfaffian systems; identifying it with the source's intended statement requires the omitted context.

Under this interpretation, fix integers \(k\geq1\) and \(1\leq q\leq k\), and let
\[
\alpha_j=\sum_{\ell=1}^k a_{j\ell}(z)\,\mathrm{d}z_\ell,
\qquad a_{j\ell}\in\mathcal{O}(\mathbb{D}^k),\quad 1\leq j\leq q.
\]
Suppose that \(\alpha_1\wedge\cdots\wedge\alpha_q\) is not identically zero and that
\[
\alpha_1\wedge\cdots\wedge\alpha_q\wedge\mathrm{d}\alpha_j=0
\qquad(1\leq j\leq q).
\]
On the locus where the forms have rank \(q\), their common kernel defines the corresponding regular foliation.

The conditional approximation question is whether, for every such system, there exist holomorphic one-forms
\[
\alpha_{j,n}=\sum_{\ell=1}^k a_{j\ell,n}(z)\,\mathrm{d}z_\ell
\qquad(1\leq j\leq q,\ n\geq1)
\]
on \(\C^k\), with all coefficients entire, respectively all coefficients polynomial, such that
\[
\alpha_{1,n}\wedge\cdots\wedge\alpha_{q,n}\not\equiv0,
\qquad
\alpha_{1,n}\wedge\cdots\wedge\alpha_{q,n}\wedge\mathrm{d}\alpha_{j,n}=0
\quad(1\leq j\leq q),
\]
and, for every compact set \(K\subset\mathbb{D}^k\),
\[
\lim_{n\to\infty}\max_{\substack{1\leq j\leq q\\1\leq\ell\leq k}}
\sup_{z\in K}|a_{j\ell,n}(z)-a_{j\ell}(z)|=0.
\]
The choice of one-forms, the displayed integrability equations, and the generic-rank hypothesis are explicit interpretive assumptions, not hypotheses recovered from the supplied excerpt. Any further requirements on singular sets or ranks must be recovered from the omitted context. The source associates polynomial approximation with foliations on \(\PP^k\), but supplies no projective extension conventions or hypotheses; the displayed coefficient approximation therefore does not by itself specify the full projective formulation.
\subsection{Short English statement}
Can every local holomorphic foliation be approximated, uniformly on compact subsets at the level of its defining forms, by entire or polynomial forms that retain the required Frobenius integrability condition? The supplied excerpt omits that condition and the defining-system hypotheses, so the explicit formulation above is conditional and requires human review.
\subsection{Sources}
\begin{itemize}
\item Local UnsolvedMath registry, permanent ID \(800690\), exact source code \texttt{catalogue entry}, title ``Holomorphic Foliation Theory --- Problem 17.'' The registry attributes the problem to Dinh and Sibony, with proposed year 2018, and labels it open. This is a supplied status label, not an independently verified mathematical or literature assessment.
\item Dinh and Sibony, \emph{Some Open Problems on Holomorphic Foliation Theory} (2018), Problem 17. Supplied source URL: \url{https://arxiv.org/abs/1807.05640}. The supplied provenance records access on 2026-10-16 and extraction by \texttt{arxiv\_latex}; these are supplied provenance claims, not a claim of independent source access or verification.
\item Supplied upstream revision: \texttt{d371ee917ac9647fbb06cd6fc529deec6f3dda82}. The complete supplied registry record contains the question but not the preceding definitions or Frobenius condition to which it refers.
\end{itemize}
\end{document}
